%\documentclass[12pt,oneside]{amsart}
\documentclass{scrartcl}

%\documentclass{article}

\usepackage{amsthm,amsmath,amssymb}
\usepackage[numbers]{natbib}

%\usepackage[colorlinks]{hyperref}
\usepackage{hyperref}
\hypersetup{
    colorlinks,%
    citecolor=black,%
    filecolor=black,%
    linkcolor=black,%
    urlcolor=black
}
\usepackage{hypernat}
\usepackage[top=2.5cm, bottom=2.5cm, left=2.8cm,
right=2.8cm]{geometry}
\usepackage{graphicx}


%\setlength{\parskip}{0pt}
%\setlength{\parsep}{0pt}
%\setlength{\headsep}{0pt}
%\setlength{\topskip}{0pt}
%\setlength{\topmargin}{0pt}
%\setlength{\topsep}{0pt}
%\setlength{\partopsep}{0pt}

\linespread{1}

%\usepackage{marvosym}

%\textwidth = 6.5 in \textheight = 9.2 in \oddsidemargin = 0.0 in
%\evensidemargin = 0.0 in \topmargin = -0.0 in \headheight = 0.0 in
%\headsep = 0.0in \parskip = 0.0in \parindent = 0.0in
%\def\baselinestretch{0.871}


\newtheorem{theorem}{Theorem}[section]
\newtheorem{proposition}[theorem]{Proposition}
\newtheorem{problem}[theorem]{Problem}
\newtheorem{fact}[theorem]{Fact}
\newtheorem{lemma}[theorem]{Lemma}
\newtheorem{defn}[theorem]{Definition}
\newtheorem{corollary}[theorem]{Corollary}
\newtheorem{remark}{Remark}[section]
\newtheorem{example}[theorem]{Example}

\newcommand{\E}{{\mathbb E}}
\newcommand{\se}{{\mathcal{E}}}
\newcommand{\R}{{\mathbb R}}
\renewcommand{\P}{{\mathbb P}}
\newcommand{\ac}{{\mathcal{A}}}
\newcommand{\A}{{\mathcal{A}}}
\newcommand{\B}{{\mathcal{B}}}
\newcommand{\C}{{\mathcal{C}}}
\newcommand{\M}{{\mathcal{M}}}
\newcommand{\Mf}{{\mathfrak{M}}}
\newcommand{\Qcal}{{\mathcal{Q}}}
\newcommand{\T}{{\mathcal{T}}}
\newcommand{\Z}{{\mathcal{Z}}}
\newcommand{\K}{{\mathcal{K}}}
\newcommand{\lc}{{\mathcal{L}}}
\newcommand{\Ps}{{\mathcal{P}}}
\newcommand{\G}{{\mathcal{G}}}
\newcommand{\pa}{{\mathring{p}}}
\newcommand{\F}{{\cal F}}
\newcommand{\samp}{{\mathcal{X}}}
\newcommand{\X}{{\mathcal{X}}}
\newcommand{\hi}{{\hat{\Phi}}}
\newcommand{\data}{{\text{data}}}
\newcommand{\relu}{{\text{ReLU}}}

\newcommand{\ridge}{{\hat{\beta}_{\text{ridge}}(\lambda)}}
\newcommand{\lasso}{{\hat{\beta}_{\text{lasso}}(\lambda)}}
\newcommand{\ridgemu}{{\hat{\mu}_t^{\text{ridge}}(\lambda)}}
\newcommand{\lassomu}{{\hat{\mu}_t^{\text{lasso}}(\lambda)}}


\newcounter{rcnt}[section]
\renewcommand{\thercnt}{(\roman{rcnt})}

\def\qt#1{\qquad\text{#1}}

\def\argmin{\mathop{\rm argmin}}
\def\argmax{\mathop{\rm argmax}}


\setlength{\parskip}{1.4 \medskipamount}
%\sloppy
%\linespread{1.3}

\begin{document}

\title{STAT 238 - Bayesian Statistics  \\
  Lecture Thirty Six} 
\subtitle{Spring 2026, UC Berkeley}

\author{Aditya Guntuboyina}


\date{27 April 2026}

\maketitle

\section{Variational Inference}

Variational Inference is a method for approximating a posterior
distribution by a simpler tractable family of distributions. It uses
optimization to select a distribution $q$ that is closest to the true
posterior in some metric. In practice, variational inference can be
much faster than MCMC, though sometimes less accurate in representing
posterior uncertainty.

The basic idea behind variational inference is very simple. Consider
the standard Bayesian setup with data given by $y$ and parameters
given by $\theta$. The prior is $f_{\theta}(\theta)$ and the
likelihood is $f_{y \mid \theta}(y)$. The posterior is then:
\begin{align*}
  f_{\theta \mid y}(\theta) = \frac{f_{\theta}(\theta) f_{y \mid
  \theta}(y)}{\int f_{\theta}(\theta) f_{y \mid \theta}(y) d\theta}. 
\end{align*}
The denominator is the marginal density $f_{y}(y)$ of $y$. The
marginal density $f_{y}(y)$is also known as the
\textbf{Evidence}. Using the marginal density for the denominator, we
get
\begin{align*}
    f_{\theta \mid y}(\theta) = \frac{f_{\theta}(\theta) f_{y \mid
  \theta}(y)}{f_{y}(y)}. 
\end{align*}
The posterior is usually intractable because of the integration
involved in the computation of the Evidence (denominator). In
variational inference, one chooses a simpler family of distributions
$\Qcal$ and then employs optimization techniques to choose a
distribution in $\Qcal$ that is as close as possible to the 
posterior $f_{\theta \mid y}(\theta)$. The distance-measure used to
measure closeness here is most commonly the following Kullback-Leibler
divergence:
\begin{align*}
  KL(q \| f_{\theta \mid y}) = \int q(\theta) \log
  \frac{q(\theta)}{f_{\theta \mid y}(\theta)} d\theta. 
\end{align*}
Recall that the KL divergence between two densities $q$ and $p$ is
given by $\int q \log (q/p)$. It is always nonnegative and it equals
zero if and only if $q = p$. It is not symmetric in general (i.e.,
$KL(q \| p) \neq KL(p \| q)$).

The optimization that we need to solve therefore is:
\begin{align}\label{viopt}
  \underset{q \in \Qcal}{\text{argmin}}~ KL(q \| f_{\theta \mid y}). 
\end{align}
The choice of this specific KL divergence is mostly for technical
convenience as it makes the resulting optimization tractable. Observe
that
\begin{align}\label{klobj}
  KL(q \| f_{\theta \mid y}) = \int q(\theta) \log
                               \frac{q(\theta)}{f_{\theta \mid
  y}(\theta)} d\theta = \int q(\theta) \log \frac{q(\theta)}{f_{y,
  \theta}(y, \theta)} 
    d\theta + \log f_{y}(y). 
\end{align}
The second term above (which involves the usually intractable
$f_y(y)$) does not depend on $q$ so the optimization \eqref{viopt} is
equivalent to: 
\begin{align*}
  \underset{q \in \Qcal}{\text{argmin}}~ \int q(\theta) \log \frac{q(\theta)}{f_{y,
  \theta}(y, \theta)} 
    d\theta. 
\end{align*}
The above objective function is called the Variational Free Energy, or
simply, Free Energy $F(q)$:
\begin{align}\label{fe}
  F(q) = \int q(\theta) \log \frac{q(\theta)}{f_{y,
  \theta}(y, \theta)} 
    d\theta = -\int q(\theta) \log \frac{f_{y, \theta}(y,
  \theta)}{q(\theta)} d\theta. 
\end{align}
So the main goal in Variational Inference is to minimize the Free
Energy. The negative of the Free Energy is called the Evidence Lower
Bound (ELBO):
\begin{align}\label{elbo}
  \text{ELBO}(q) = -F(q) = \int q(\theta) \log \frac{f_{y, \theta}(y,
  \theta)}{q(\theta)} d\theta. 
\end{align}
So the main goal in Variational Inference can also be said to be the
maximization of the ELBO. The name ELBO comes from the fact that
$\text{ELBO}(q)$ always satisfies:
\begin{align}\label{ev_lb}
  \text{ELBO}(q) = \int q(\theta) \log \frac{f_{y, \theta}(y,
  \theta)}{q(\theta)} d\theta \leq \log f_y(y). 
\end{align}
This is because of the expression \eqref{klobj} which says that the
difference between $\log f_{y}(y)$ and $\text{ELBO}(q)$ equals the
Kullback Leibler divergence $KL(q \| f_{\theta \mid y})$ which is
nonnegative. Because $f_y(y)$ is called the evidence and
$\text{ELBO}(q)$ gives a lower bound for the logarithm of the
evidence, it is given the name Evidence Lower Bound. It is also easy
to see that if we maximize $\text{ELBO}(q)$ over \textit{all}
probability densities $q$, then we obtain equality in \eqref{ev_lb}:
\begin{align}\label{sup_elbo}
  \sup_q \text{ELBO}(q) = \log f_{y}(y). 
\end{align}
The above follows from \eqref{klobj} because
\begin{align*}
  \text{ELBO}(q) = \log f_{y}(y) - KL(q \| f_{\theta \mid y}). 
\end{align*}

Because $f_{y, \theta}(y, \theta) = f_{\theta}(\theta) f_{y \mid
  \theta}(y)$, the ELBO has the following alternative expression:
\begin{align}
  \text{ELBO}(q) &= -F(q) \nonumber \\ &= \int q(\theta) \log \frac{f_{y, \theta}(y,
  \theta)}{q(\theta)} d\theta \nonumber \\ &= \int q(\theta) \log \frac{f_{y\mid
                                   \theta}(y)
                                   f_{\theta}(\theta)}{q(\theta)}
                                   d\theta \nonumber \\
  &= \int q(\theta) \log f_{y \mid \theta}(y) d\theta - \int q(\theta)
    \log \frac{q(\theta)}{f_{\theta}(\theta)} d\theta \nonumber \\
  &= \int q(\theta) \log f_{y \mid \theta}(y) d\theta - KL(q \|
    f_{\theta}) \label{elbo_alt}
\end{align}
In other words, $\text{ELBO}(q)$ represents the average likelihood
(with respect to $q$) minus the KL divergence between $q$ and the
prior.

Usually, the main tasks in Variational Inference are (a) to choose a
class of distributions $\Qcal$, and (b) to maximize $\text{ELBO}(q)$
(or, equivalently, to minimize $F(q)$) over $q \in \Qcal$. We will see
how to do this via some examples.

\section{Bayesian Logistic Regression}
Consider again the logistic regression setting where we observe data
$(x_i, y_i), i = 1, \dots, n$ with $x_i \in \R^d$ and $y_i \in \{0,
1\}$. The likelihood model is:
\begin{align*}
  y_i \mid x_i, \beta \sim \text{Bernoulli}\left(\frac{e^{x_i^T
  \beta}}{1 + e^{x_i^T \beta}} \right)
\end{align*}
and we take the improper uniform prior for $\beta$. We then have:
\begin{align*}
  f_{y, \beta}(y, \beta) &= \prod_{i=1}^n \left(\frac{e^{x_i^T
  \beta}}{1 + e^{x_i^T \beta}} \right)^{y_i} \left(1 - \frac{e^{x_i^T
                           \beta}}{1 + e^{x_i^T \beta}} \right)^{1-y_i} \\
  &= \prod_{i=1}^n \frac{\exp(y_i x_i^T \beta)}{1 + \exp(x_i^T
    \beta)} = \exp(\ell(\beta))
\end{align*}
where $\ell(\beta)$ is the log-likelihood is
\begin{align}\label{loglik}
  \ell(\beta) = \sum_{i=1}^n \left(y_i x_i^T \beta - \log \left(1 +
  \exp(x_i^T \beta) \right) \right).
\end{align}
Note that, because the prior is one, $f_{y, \beta}(y, \beta)$ equals
the likelihood $f_{y \mid \beta}(y)$. We can thus calculate the ELBO as:
\begin{align*}
  \text{ELBO}(q) &= \int q(\beta) \log \frac{f_{y, \beta}(y, \beta)}{q(\beta)} d\beta \\
  &= \int q(\beta) \ell(\beta) d\beta - \int q(\beta) \log q(\beta) d\beta.
\end{align*}
The quantity $-\int q \log q$ is known as the entropy $H(q)$ of
$q$. We thus have
\begin{align*}
  \text{ELBO}(q) = \int q(\beta) \ell(\beta) d\beta - H(q). 
\end{align*}
Let us now take $\Qcal$ to be the class of all normal distributions $N(m,
\Sigma)$ as $m \in \R^p$ and $\Sigma$ varies over the class of all $p
\times p$ positive definite matrices $\Sigma$. The entropy of the
multivariate normal distribution (see
\url{https://en.wikipedia.org/wiki/Multivariate_normal_distribution})
is:
\begin{align*}
  H(N(m, \Sigma)) = \frac{p}{2} \left(1 + \log (2 \pi) \right) +
  \frac{1}{2} \log |\Sigma|
\end{align*}
where $|\Sigma|$ is the determinant of $\Sigma$. We thus have:
\begin{align*}
  \text{ELBO}(m, \Sigma) = \E_{\beta \sim N(m, \Sigma)} \ell(\beta) +
  \frac{1}{2} \log |\Sigma| + \frac{p}{2} \left(1 + \log(2 \pi)
  \right). 
\end{align*}
This can be maximized over $m$ and $\Sigma$ to obtain the posterior
approximation $N(\hat{m}, \hat{\Sigma})$. Here $m$ can be any vector
over $\R^p$ but $\Sigma$ is constrained to be positive
definite. Constrained optimization is usually difficult (for
gradient-based methods) so we convert this to unconstrained
optimization by taking
\begin{align*}
  \Sigma = L L^T
\end{align*}
where $L$ is a $p \times p$ lower-triangular matrix with non-zero
diagonal entries. We can also, without loss of generality, assume that
$L$ has strictly positive diagonal entries (given any $L$, we can
replace it by $L D$ where $D$ is a diagonal matrix with entries $+1$
if the corresponding entry of $L$ is positive and $-1$
otherwise). In code, we can represent the diagonal entries of $L$ by
$\exp(a_1), \dots, \exp(a_p)$. With this parametrization, we have
\begin{align*}
|\Sigma| = |L L^T| = |L|^2 = \prod_{j=1}^p L_{jj}^2, 
\end{align*}
and thus
\begin{align*}
  \text{ELBO}(m, L) = \E_{\beta \sim N(m, L L^T)} \ell(\beta) +
  \sum_{j=1}^p \log L_{jj} + \frac{p}{2} \left(1 + \log(2 \pi)
  \right). 
\end{align*}
In order to maximize the above function, we need to make the term
$\E_{\beta \sim N(m, L L^T)} \ell(\beta)$ more explicit. For this, we
use the \textit{reparametrization} trick (see
\url{https://en.wikipedia.org/wiki/Reparameterization_trick}) and
write
\begin{align*}
  \beta = m + L z ~~ \text{ where $z \sim N(0, I_p)$}. 
\end{align*}
This gives
\begin{align*}
  \text{ELBO}(m, L) = \E_{z \sim N(0, I_p)} \ell(m + L z) +
  \sum_{j=1}^p \log L_{jj} + \frac{p}{2} \left(1 + \log(2 \pi)
  \right).  
\end{align*}
The expectation above can be approximated by Monte Carlo. Specifically
generate
\begin{align*}
  z^{(1)}, \dots, z^{(S)} \overset{\text{i.i.d}}{\sim} N(0, I_p)
\end{align*}
and use the approximation:
\begin{align*}
  \text{ELBO}(m, L) \approx \frac{1}{S} \sum_{s=1}^S \ell(m + L
  z^{(s)}) +
  \sum_{j=1}^p \log L_{jj} + \frac{p}{2} \left(1 + \log(2 \pi)
  \right).  
\end{align*}
This can be maximized using standard gradient-based optimization
software such as PyTorch. 



%\begin{thebibliography}{99}

%  \bibitem[Neal(2011)]{neal2011mcmc}
%Radford M. Neal.
%\newblock MCMC using Hamiltonian dynamics.
%\newblock In \emph{Handbook of Markov Chain Monte Carlo}, pages 47--95.
%\newblock Chapman and Hall/CRC, 2011.

%\bibitem[Girolami and Calderhead(2011)]{girolami2011riemann}
%Mark Girolami and Ben Calderhead.
%\newblock Riemann manifold Langevin and Hamiltonian Monte Carlo methods.
%\newblock \emph{Journal of the Royal Statistical Society: Series B (Statistical Methodology)}, 73(2):123--214, 2011.

%  \bibitem[Lai et~al.(2025)Lai, Song, Kim, Mitsufuji, and Ermon]{lai2025principles}
%Lai, C.-H., Song, Y., Kim, D., Mitsufuji, Y., and Ermon, S. (2025).
%\newblock The principles of diffusion models.
%\newblock \emph{arXiv preprint arXiv:2510.21890}.

%\bibitem[Glatt-Holtz et~al.(2024)]{glatt2024sacred}
%Nathan E. Glatt-Holtz, Andrew J. Holbrook, Justin A. Krometis, Cecilia F. Mondaini, and Ami Sheth.
%\newblock Sacred and Profane: from the Involutive Theory of MCMC to Helpful Hamiltonian Hacks.
%\newblock \emph{arXiv preprint arXiv:2410.17398}, 2024.

%\bibitem[Roberts and Rosenthal(2004)]{RobertsRosenthal2004}
%Roberts, G. O. and Rosenthal, J. S. (2004).
%General state space Markov chains and MCMC algorithms.
%\textit{Probability Surveys}, \textbf{1}, 20--71.

%\bibitem[Meyn and Tweedie(2009)]{MeynTweedie2009}
%Meyn, S. P. and Tweedie, R. L. (2009).
%\textit{Markov Chains and Stochastic Stability}, 2nd ed.
%Cambridge University Press.


  
%\bibitem[Chib and Greenberg(1995)]{chib1995}
%Chib, S. and Greenberg, E. (1995).
%Understanding the Metropolis--Hastings algorithm.
%\textit{The American Statistician}, \textbf{49}, 327--335.


%\bibitem[MacKay and Peto(1995)]{MacKay1995HierarchicalDirichlet}
%D.~J.~C. MacKay and L.~C. Peto,
%\newblock ``A hierarchical Dirichlet language model,''
%\newblock \emph{Natural Language Engineering}, vol.~1,
%no.~3, pp.~289--308, 1995.

%\bibitem[Rubin(1981)]{Rubin1981BayesianBootstrap}
%D.~B. Rubin,
%\newblock ``The Bayesian bootstrap,''
%\newblock \emph{The Annals of Statistics}, vol.~9,
%no.~1, pp.~130--134, 1981.
  
%\bibitem[Fisher(1934)]{Fisher1934}
%R.~A. Fisher,
%\newblock ``Probability, likelihood and quantity of information in the logic of
%uncertain inference,''
%\newblock \emph{Proceedings of the Royal Society of London. Series~A,
%Containing Papers of a Mathematical and Physical Character}, vol.~146,
%no.~856, pp.~1--8, 1934.

%\bibitem[Efron(2010)]{Efron}
%Efron, B. (2010)
%\textit{Large-scale inference: empirical Bayes methods for estimation,
%testing, and prediction}. Cambridge University Press, 2010.

%\bibitem[Shumway and Stoffer(2017)]{ShumwayStoffer}
%Shumway, R. H. and Stoffer, D. S. (2017)
%\textit{Time Series Analysis and Its Applications: With R
%  Examples}. Springer, 4th edition. 
  
%     \bibitem[Jaynes(2003)]{Jaynes} Jaynes, E. T, (2003)
%  \textit{Probability theory: the logic of science}. Cambridge
%  University Press, 2003.
 
%  \bibitem[Sinai(1992)]{Sinai} Sinai, Y. G, (1992)
%  \textit{Probability theory: an introductory course}. Springer
%  Verlag, 1992.  

%\bibitem[{deGroot(1986)}]{DeGrootBlackwell}DeGroot, M. H. (1986)
%       A conversation with David Blackwell.
%\newblock \emph{Statistical Science}, \textbf{1}, 40--53.

%         \bibitem[Mosteller(1987)]{Mosteller} Mosteller, F, (1987)
%  \textit{Fifty challenging problems in probability with
%    solutions}. Courier Corporation, 1987.

%    \bibitem[MacKay(2003)]{MacKay} MacKay, D. J. C., (2003)
%  \textit{Information theory, inference, and learning
%  algorithms}. Cambridge University Press, 2003.

%\bibitem[{Lindley(2013)}]{Lindley}Lindley, D. V. (2013)
%   \textit{Understanding Uncertainty}. John Wiley and sons, 2013.


%\end{thebibliography}







\end{document}

%%% Local Variables:
%%% mode: latex
%%% TeX-master: t
%%% End:
